\chapter{Sequential Model Based Optimization}

There has been a lot of work done over the last few years in the space of Bayesian Optimization. One primary advance has been the formalization of sequential model based optimization (SMBO) as a general framework for Bayesian Optimization. The general idea of \green{SMBO methods} is to divide the optimization problem into two stages, first fitting a model from a specific class that predicts the performance on the target function and then gathering additional data to further refine the model that is used as a surrogate for the actual function we wish to evaluate.

\green{Hyperparameter tuning is worth pausing on before the algorithm is stated, because it is not a new problem sitting alongside those of the previous chapter but a direct instance of them. A learning algorithm's hyperparameters span a search space $\mathcal{X}$; training and validating the model at one configuration $x \in \mathcal{X}$ is one evaluation of an expensive black-box objective $f$; and the validation score is the quantity we want to optimize. Nothing in the Bayesian optimization machinery changes under this identification. It is also the setting where the framework earns its keep, since a single evaluation can cost hours of compute, and it is the application that motivated much of the field's recent development. For these reasons it serves as the running example of this chapter and the benchmark family of the experimental chapter.}

As detailed in \green{\cite{hutter2011sequential}\footnote{\green{For an accessible informal exposition of the same loop, see \cite{conceptualBO}.}}}, the two steps can be broken down into this pseudo-algorithm for the problem of finding the best hyperparameters for a generic machine learning model: 

\begin{enumerate}
	\item Build a surrogate probability model for the target function
	\item Find the hyperparameters that perform the best on the surrogate
	\item Apply these hyperparameters to the true objective function\footnote{This is the part that is expensive to evaluate, ideally, we want to minimize the amount of times that we actually run the true objective function.}
	\item Update the surrogate model incorporating the new results
	\item Repeat steps 2-4 until a \green{stopping condition is met} (maximum number of iterations, time spent, or others such as cost).
\end{enumerate}

\green{To see the loop in motion, consider a generic instance of it. Suppose the search space is a bounded interval, the budget allows ten evaluations, and we spend the first three on an initial design, evaluating $f$ at three spread-out points and fitting the surrogate to the resulting pairs (step 1). The surrogate now carries a prediction and an uncertainty at every point of the interval. The selection criterion scores candidates by combining the two, and suppose it peaks at a point $x_4$ lying between the two best observations, where the prediction is good and the uncertainty still moderate (step 2). We pay for one true evaluation at $x_4$ (step 3), observe a value slightly better than anything seen so far, and refit the surrogate on the enlarged dataset (step 4).}

\green{On the next pass the picture has changed. Around $x_4$ the surrogate is now confident, so the criterion no longer favours that neighbourhood as strongly, and its maximum moves to a stretch of the interval that has never been probed, where uncertainty is high (step 2 again). The true function is evaluated there (step 3), the result turns out to be mediocre, the surrogate absorbs that information (step 4), and attention shifts back toward the earlier promising region. The loop continues in this way, alternating between sharpening the picture where good values live and ruling out unexplored territory, until the ten evaluations are spent (step 5). Every surrogate model surveyed in the next chapter slots into exactly this cycle; what changes between methods is only how the predictions and uncertainties that drive steps 1 and 4 are produced.}

While all SMBO algorithms follow these basic steps, the true power of \green{these algorithms} is in the flexibility \green{they provide} when selecting the surrogate model and when selecting how the new results can be incorporated into the model itself (steps 3 and 4). This flexibility is the main point of the thesis, so it is worth dwelling on. Nothing in the SMBO template requires the surrogate to be a Gaussian process; the template asks only for a model that predicts the objective and carries some notion of uncertainty about its predictions, and any model meeting that description can be slotted in. This matters because the surrogate is not a neutral choice: it silently determines which problems the optimizer can even represent. A Gaussian process, the usual default, builds its predictions from a kernel that measures the numeric distance between inputs, and so it presumes a continuous search space. Swap the surrogate for one that does not make that presumption, and the same SMBO loop suddenly applies to problems the GP cannot natively express, most notably problems with categorical inputs. That substitution is exactly what this thesis carries out: it keeps the SMBO framework fixed and replaces the surrogate with a Bayesian Adaptive Spline Surface, whose treatment of categorical variables is native rather than bolted on. The acquisition step (the ``how to incorporate new results'' half of the flexibility) is held fixed across the methods we compare, precisely so that the effect of changing the surrogate can be isolated.

As is clear by the title of this work, the main contribution is the implementation of the BASS model as a surrogate model under this framework, but there have been a large amount of different methods and models used as surrogate models as well as a large amount of ways of incorporating the results into the surrogate, known as the selection function or evaluation criteria depending on the literature being read. The next sections detail the other models that are commonly used as surrogate models, and where these models got their roots.

